\scanpage{134}
\mspage{66}
termes, il faut introduire, \deleted{\unsure{des cas}}
pour un morphisme
\[
(218)
\]
\diagram{\begin{tikzcd}[ampersand replacement=\&]\Delta_r\arrow[r,no head,"\alpha"]\&\Delta_{r'}\end{tikzcd}},
un \uline{morphisme correspondant} de groupoïdes
\[
(219)
\]
\diagram{\fbox{\begin{tikzcd}[ampersand replacement=\&]\widetilde\Pi_{\Delta_{r'}}\arrow[r,"\widetilde\alpha^*"]\&\Pi_{\Delta_r}\end{tikzcd}}}
qu'on appellera l'\uline{hom de générisation}
associé à $\alpha$.
\margin{\figure{\begin{tikzpicture}\draw (0,0)--(0,1);\draw (.05,0)--(.05,1);\end{tikzpicture}}}
Comme $\alpha$ se factorise
canoniquement en un hom surjectif
suivi d'un hom injectif de simplexes,
on est ramené à examiner séparément
les deux cas. Mais si $\alpha$ est surjectif
(« hom de dégénérescence »), \deleted{l'homomorphisme}
\inserted{on a $\alpha(0)=0$}
donc on est dans le cas déjà envisagé --
\deleted{\unsure{correspondant}}\inserted{d'ailleurs} dans ce cas
\diagram{\begin{tikzcd}[ampersand replacement=\&]D^*_{\Delta_{r'}}\arrow[r]\&D^*_{\Delta_r}\end{tikzcd}}
\unsure{se}
décompose en des morphismes
\[
(220)
\]
\diagram{\begin{tikzcd}[ampersand replacement=\&]
D^*_{d'_0,\ldots,d'_{r'}}\arrow[r]\&D^*_{\unsure{d_0|d'_0},\ldots,\unsure{d_r|d'_r}}
\end{tikzcd}},
$(d_0,\ldots,d_r)$ est déduite de $(d'_0,\ldots,\unsure{d'_r|d'_{r'}})$
en \unsure{réécrivant} cette suite, en répétant certains
de ses éléments -- et \deleted{les} alors (220) est
un iso -- donc \uline{si $\alpha$ est un hom}
\margin{On voit que la considération des $\widetilde\Pi_{d_0- d_r}$, $\Pi_{d_0,-d_r}$ pour $d_0,\ldots,d_r$ \unsure{non} \uline{strictement} décroissants, et celle de $\alpha^*$, pour $\alpha$ \unsure{non} \uline{stricts} croissants, \unsure{est de la poudre aux yeux}, \unsure{simplement} !.}
\margin{133}
